\section{Methodology}\label{sec:method}
Figure~\ref{fig:overview-proposed-method} summarizes the proposed pipeline for
constructing and exposing a movie knowledge graph. Movie data are collected from TMDB and stored in CSV
files. These tabular records are transformed into Turtle RDF using stable
resource URIs. The OWL ontology then adds class, property, and restriction
semantics and supports OWL-RL inference. Movie entities are linked to DBpedia
and Wikidata using identifiers, titles, release years, and related matching
criteria. Apache Jena Fuseki provides a browser-accessible SPARQL endpoint for
querying the graph. The optional static web application publishes generated
resource pages and ontology views for human browsing.

Here, we describe the first step of the workflow
(i.e., defining the ontology for the movie domain) and specify the evaluation protocol,
which assesses the resulting knowledge graph by executing representative SPARQL queries. The
remaining stages are described in their corresponding sections.

\paragraph{Movie ontology}

We construct the movie ontology by reusing Schema.org classes and properties
to represent movies and their related entities. Table~\ref{tab:schema-vocabulary}
lists all Schema.org terms used in the ontology~\citep{schemaorg2025}. The classes model movies,
people, roles, descriptive entities, organizations, and ratings, while the
properties describe movie metadata and relationships, including titles,
identifiers, release dates, duration, genres, languages, countries, production
companies, ratings, directors, actors, and acting roles. This vocabulary
provides a shared semantic model while allowing the data to be serialized as
RDF and constrained or extended with OWL.

\paragraph{Evaluation protocol}

The evaluation protocol queries the graph using SPARQL. The query set covers
basic retrieval, such as listing movies, people, and directors, as well as
multi-condition questions involving ratings, genres, release years, and cast
members. Additional queries test the value of OWL-RL reasoning by asking for
inferred relationships, such as a person's directed movies or a production
company's country. Finally, federated queries retrieve information about a
movie from DBpedia or Wikidata through its established links. Queries are
executed against the local graph, with reasoning enabled when required, and
against the corresponding remote endpoint for external-source questions.

\begin{table}[t]
    \centering
    \caption{Schema.org classes and properties reused in the movie ontology.}
    \label{tab:schema-vocabulary}
    \small
    \begin{tabular}{p{0.22\linewidth}p{0.70\linewidth}}
        \toprule
        \textbf{Category} & \textbf{Schema.org terms and use} \\
        \midrule
        Classes &
        \texttt{Movie} (movie); \texttt{CreativeWork} (movie superclass);
        \texttt{Person} (actor, director, or crew member);
        \texttt{Role} (reified acting role); \texttt{Genre} (movie genre);
        \texttt{Country} (movie or company country);
        \texttt{Language} (movie language);
        \texttt{Organization} (production company);
        \texttt{AggregateRating} (movie rating);
        \texttt{Rating}, \texttt{Thing}, and \texttt{Intangible}
        (reused superclasses). \\
        \midrule
        Object properties &
        \texttt{actor} (movie or role to actor);
        \texttt{aggregateRating} (movie to rating);
        \texttt{contributor} (movie to non-director crew member);
        \texttt{countryOfOrigin} (movie to country);
        \texttt{director} (movie to director);
        \texttt{genre} (movie to genre);
        \texttt{inLanguage} (movie to language);
        \texttt{location} (company to country);
        \texttt{productionCompany} (movie to company);
        \texttt{sameAs} (movie to an equivalent external resource). \\
        \midrule
        Datatype properties &
        \texttt{abstract} (movie synopsis);
        \texttt{characterName} (character portrayed in a role);
        \texttt{datePublished} (release date);
        \texttt{duration} (runtime);
        \texttt{identifier} (TMDB or IMDb identifier);
        \texttt{name} (entity name);
        \texttt{position} (cast order);
        \texttt{ratingCount} (number of votes);
        \texttt{ratingValue} (rating value);
        \texttt{roleName} (crew job title). \\
        \bottomrule
    \end{tabular}
\end{table}